\documentclass[10pt]{article}
\usepackage[T1]{fontenc}
\usepackage{charter}
\usepackage[letterpaper,margin=0.55in,top=0.45in,bottom=0.45in]{geometry}
\usepackage{titlesec}
\usepackage{enumitem}
\usepackage{xcolor}
\usepackage{tabularx}
\usepackage[hidelinks]{hyperref}
\usepackage{microtype}

\definecolor{accent}{RGB}{20,50,100}
\pagestyle{empty}
\setlength{\parindent}{0pt}
\linespread{1.03}

\titleformat{\section}{\color{accent}\bfseries\large\scshape}{}{0pt}{}[\vspace{-1.2ex}\color{accent}\rule{\linewidth}{0.6pt}]
\titlespacing*{\section}{0pt}{0.9ex}{0.2ex}

\setlist[itemize]{leftmargin=1.2em,itemsep=0.15ex,topsep=0.2ex,parsep=0pt,label={\small\textbullet}}

% \entry{Left bold}{Right}{Left italic}{Right italic}
\newcommand{\entry}[4]{%
  \begin{tabularx}{\linewidth}{@{}X r@{}}
    \textbf{#1} & #2 \\
    \textit{#3} & \textit{#4} \\
  \end{tabularx}\vspace{-0.2ex}}

\begin{document}

\begin{center}
  {\LARGE\bfseries Wenfeng Wang}\\[0.6ex]
  Philadelphia, PA \;$\cdot$\; \href{mailto:fw93810@engineering.upenn.edu}{fw93810@engineering.upenn.edu} \;$\cdot$\; \href{https://fw93810.github.io}{fw93810.github.io}
\end{center}
\vspace{-1ex}

\section{Education}
\entry{University of Pennsylvania}{Philadelphia, PA}{M.S.E.\ in Electrical Engineering \quad GPA: 3.72/4.00}{Aug 2025 -- May 2027 (expected)}
\begin{itemize}
  \item Coursework: Machine Learning, Graph Neural Networks, Statistics for Data Science, Digital Integrated Circuits \& VLSI, Artificial Intelligence*, Linear Systems Theory* \hfill (*in progress)
\end{itemize}
\vspace{0.4ex}
\entry{Jilin University}{Changchun, China}{B.S.\ in Physics \quad GPA: 3.80/4.00}{Sep 2020 -- Jul 2024}
\begin{itemize}
  \item Coursework: Electrodynamics, Quantum Mechanics, Thermal Physics, Mathematical Analysis, Linear Algebra, Probability
\end{itemize}

\section{Publications}
\begin{enumerate}[leftmargin=1.6em,itemsep=0.3ex,topsep=0.2ex,label={[\arabic*]}]
  \item \textbf{Wenfeng Wang}, Jianrui Shi, Mingang Guo, Jiayuan Ma, Yuqi Lu, Jiajun Sun, Guojiang Zhao, Jing Li.
    Trace-to-Architecture: Capability Bottleneck Diagnosis for Harness Revision.
    \textit{arXiv preprint} \href{https://arxiv.org/abs/2610.16725}{\underline{arXiv:2610.16725}}, 2026. Under review at ACL Rolling Review.
  \item \textbf{Wenfeng Wang}.
    The Design and Optimization for 4-Bit Absolute Value Comparator in Voice Signal Processing.
    Accepted for conference proceedings, 2023 (not yet published).
  \item \textbf{Wenfeng Wang}.
    The Introduction of Design and Research about PeLED.
    Accepted for conference proceedings, 2023 (not yet published).
\end{enumerate}

\section{Research Experience}
\entry{Trace-to-Architecture: Capability Bottleneck Diagnosis for Harness Revision}{2026 -- present}{First author}{}
\begin{itemize}
  \item Proposed T2A, a framework that uses execution trajectories to decide what to change in an LLM agent's harness when it fails, and when those changes should apply.
  \item Built tooling that groups trajectories by observable features, letting a development agent diagnose failures with evidence from both successful and failed tasks.
  \item Designed validation that checks not only whether target tasks improve, but also whether selected successful tasks stay successful and whether the edits are actually triggered.
  \item Across ALFWorld, ScienceWorld, and AppWorld, achieved the highest aggregate score for all three task agents, exceeding the strongest baselines by 6.2--30.3 points; integrating T2A's diagnosis and validation into Meta-Harness raised its mean ScienceWorld score by 21.6 points.
\end{itemize}
\vspace{0.4ex}
\entry{Institute of Microelectronics, Chinese Academy of Sciences}{Beijing, China}{Research Intern}{Jul 2024 -- Aug 2025}
\begin{itemize}
  \item Studied automatic evaluation of text-to-image generation (CLIPScore, VQAScore, BLIP-2), focusing on how well metrics capture spatial relations between objects.
  \item Explored a 3D-aware evaluation dimension: recovered scene structure with monocular 3D reconstruction and fed text, image, and 3D tokens jointly into a language model to assess spatial orientation and agent--patient interactions.
  \item Reproduced and deployed vision and generation models (VGG16/19, YOLOv5, Omost) in PyTorch.
\end{itemize}
\vspace{0.4ex}
\entry{Introduction to VLSI Research Program}{Remote}{Supervised by Prof.\ Dejan Markovi\'c, UCLA}{Jul 2023 -- Aug 2023}
\begin{itemize}
  \item Led a student team; designed, built, and tested most of the circuits; wrote the final report and gave the oral presentation.
\end{itemize}
\vspace{0.4ex}
\entry{Nanoelectronics, Optoelectronics and AI Research Seminar}{Remote}{Supervised by Prof.\ Deep Jariwala, University of Pennsylvania}{Feb 2023 -- Apr 2023}
\begin{itemize}
  \item Studied nanoscale devices, nanofabrication, and low-dimensional semiconductors for AI hardware; led weekly team assignments, gave the oral presentation, and wrote the PeLED design sections that became publication [3].
\end{itemize}

\section{Awards}
\begin{tabularx}{\linewidth}{@{}X r@{}}
  Second Prize Scholarship, Jilin University & 2022 \\
  Yan Jici Scholarship, Institute of Physics, Chinese Academy of Sciences & 2021 \\
  Third Prize Scholarship, Jilin University & 2020 \\
\end{tabularx}

\section{Skills}
\textbf{Programming:} Python, PyTorch, C, LaTeX \qquad \textbf{Languages:} Chinese (native), English (fluent)

\end{document}
